\documentclass{article}

\usepackage{fancyhdr}
\usepackage{extramarks}
\usepackage{amsmath}
\usepackage{amsthm}
\usepackage{amsfonts}
\usepackage{tikz}
\usepackage[plain]{algorithm}
\usepackage{algpseudocode}

\usetikzlibrary{automata,positioning}

\newcommand{\overbar}[1]{\mkern 1.5mu\overline{\mkern-1.5mu#1\mkern-1.5mu}\mkern 1.5mu}

%
% Basic Document Settings
%

\topmargin=-0.45in
\evensidemargin=0in
\oddsidemargin=0in
\textwidth=6.5in
\textheight=9.0in
\headsep=0.25in

\linespread{1.1}

\pagestyle{fancy}
\lhead{Yousef Alaa Awad}
\chead{\hmwkClass\: \hmwkTitle}
\rhead{\firstxmark}
\lfoot{\lastxmark}
\cfoot{\thepage}

\renewcommand\headrulewidth{0.4pt}
\renewcommand\footrulewidth{0.4pt}

\setlength\parindent{0pt}

%
% Create Problem Sections
%

\setcounter{secnumdepth}{0}
\newcounter{partCounter}
\newcounter{homeworkProblemCounter}
\setcounter{homeworkProblemCounter}{1}

\newcommand{\hmwkTitle}{Homework\ \#2}
\newcommand{\hmwkDueDate}{February 6, 2026}
\newcommand{\hmwkClass}{Advanced Computer Architecture}

%
% Title Page
%

\title{
    \vspace{2in}
    \textmd{\textbf{\hmwkClass:\ \hmwkTitle}}\\
    \normalsize\vspace{0.1in}
    \vspace{3in}
}

\author{Yousef Alaa Awad}

% Problems start here
\begin{document}

\maketitle
\pagebreak

\section{1: Performance}
Suppose that we have two machines A and B. The following table shows the execution times for programs that make up a benchmark suite for machine A, machine B, and a reference machine X.
\begin{table}[h]
  \centering
  \begin{tabular}{|c|c|c|c|}
    \hline
    Program & Exec time of X & Exec time of A & Exec time of B \\ \hline
    P1      & 120            & 40             & 20             \\ \hline
    P2      & 300            & 100            & 150            \\ \hline
    P3      & 50             & 5              & 10             \\ \hline
    P4      & 150            & 25             & 30             \\ \hline
    P5      & 90             & 10             & 15             \\ \hline
  \end{tabular}
\end{table}

\subsection{Calculate the overall speedup of machine A versus reference machine X, using geometric mean.}
First we have to find the relative speedup from x of every program's execution time to machine A and B.
\begin{table}[h]
  \centering
  \begin{tabular}{|c|c|c|c|}
    \hline
    Program & Exec time of X & Exec time of A & Exec time of B \\ \hline
    P1      & 1x             & 3x             & 6x             \\ \hline
    P2      & 1x             & 3x             & 2x             \\ \hline
    P3      & 1x             & 10x            & 5x             \\ \hline
    P4      & 1x             & 6x             & 5x             \\ \hline
    P5      & 1x             & 9x             & 6x             \\ \hline
  \end{tabular}
\end{table}
Now, given the Geometric Mean equation below...
$$ Geometric\ Mean = \sqrt[n]{\left( \prod_{i=1}^{n} x_i \right)} $$
We can calculate the Geometric Mean for machine A, and machine B:
\begin{align*}
  GMean_{A} = \sqrt[5]{\left( 3*3*10*6*9 \right)} = \sqrt[5]{\left( 4860 \right)} &= 5.4617 \\
  GMean_{B} = \sqrt[5]{\left( 6*2*5*5*6 \right)} = \sqrt[5]{\left( 1800 \right)}  &= 4.4777
\end{align*}
Therefore, the overall speedup of Machine A vs the reference, X, is 5.4617 times faster, and the overall speedup of Machine B (despite not being asked formally) vs the reference, X, is 4.4777 times faster.


\pagebreak
\section{2: Reliability}
Consider a VLSI Microprocessor has a failure rate of 400 FITs. If a VLSI microprocessor has $10^5$ transistors, and the transistors have equal failure rates and transistors fail independently from one another.

\subsection{What is the MTTF of a single transistor?}
Now, we know that...
$$ \lambda_{transistors} = \frac{\lambda_{total}}{n_{transistors}} = \frac{400 \text{ FITs}}{10^5 \text{ transistors}} = 0.004\ FITs $$
But now we need to find the "failures per hour" by doing the following:
$$ \text{failures per hour} = \lambda_{total} * 10^{-9} = 4*10^{-12} \frac{failures}{hour} $$
Now, we can finally calculate MTTF:
$$ MTTF = \frac{1}{failures\ per\ hours} = \frac{1}{4*10^{-12}} = 2.5*10^{11} \text{ hours} $$


\section{3: Amdahl's Law}
Your pipeline has one very slow stage that consumes 60\% of the time to process an instruction. You discover that you can speed it up by a factor of 3

\subsection{What is the improvement in the latency?}
Now, we will have to use this equation for speedup...
$$ Speedup = \frac{time_{old}}{time_{new}} = \frac{1}{(1-f)+\frac{f}{s}} $$
Alongside this, we were given the following from the problem proper:
$$ f = 60\% = 0.6 $$
$$ s = 3 $$
Therefore, we simply just plug those values into the speedup equation above!
$$ Speedup = = \frac{1}{(1-f)+\frac{f}{s}} = \frac{1}{(1-0.6)+\frac{0.6}{3}} = \frac{1}{0.4 + 0.2} = \frac{1}{0.6} = 1.\overbar{66} $$
Therefore the total improvement in latency is $1.\overbar{66}$ times.


\end{document}
